% sec:geo-termination — Termination Conditions
\section{Termination Conditions}
\label{sec:geo-termination}

Every ray must end somewhere: absorbed by the black hole, escaped to the
celestial sphere, intercepted by the accretion disc, or abandoned because
the computation has exceeded its budget.  Without explicit stopping
criteria, the integrator will advance a geodesic indefinitely.  Each
outcome produces a different pixel in the final image --- shadow, sky-map
lookup, disc emission, or artefact --- so the termination criterion is
not merely an implementation detail but a direct determinant of the
rendered output.
\coderef{Geodesic.Termination}{checkTermination}

\paragraph{The termination type.}

The codebase encodes every possible stopping reason in a single algebraic
data type:
%
\begin{lstlisting}[language=Haskell]
data TerminationReason
  = HitHorizon !Double
  | HitCelestialSphere !Double
  | MaxAffineReached !Double
  | StepSizeTooSmall !Double
  | HitAccretionDisk !Double
  | OutOfGrid !Double
\end{lstlisting}
%
Each constructor carries a diagnostic payload --- typically the
coordinate radius at which the event occurred.  The rendering pipeline
(\cref{ch:ray-tracing}) pattern-matches on this type to determine the
pixel colour.  A \texttt{HitHorizon} result produces a shadow pixel.
\texttt{HitCelestialSphere} triggers a sky-map lookup at the final
ray direction.  \texttt{HitAccretionDisk} invokes the disc emission
model at the crossing radius.  The remaining constructors signal
computational limits rather than physical events.
\coderef{Geodesic.Types}{TerminationReason}

\paragraph{Horizon crossing.}

A photon that crosses the event horizon cannot return to a distant
observer --- the one-way membrane property guarantees this regardless of
the interior geometry --- so continuing the integration past the horizon
wastes computation without contributing to the image.  The codebase
terminates a ray when its coordinate radius falls below a threshold
%
\begin{equation}
  r < r_+ + \delta \,,
  \label{eq:horizon-buffer}
\end{equation}
%
where $r_+$ is the outer horizon radius and $\delta$ is a small buffer.
The default is $\delta = 0.01\,M$.  The buffer avoids consuming
integration steps in the strong-field region just inside the horizon
where the step size shrinks rapidly (\cref{sec:geo-raytracing}), and
provides a safety margin against floating-point roundoff in the radius
computation.

In Kerr--Schild coordinates, the metric is regular at the horizon ---
there is no coordinate singularity to contend with, unlike
Boyer--Lindquist coordinates where $\Delta = 0$ makes the metric
components diverge (\cref{sec:kerr-bl}).  The integrator could in
principle continue through $r_+$ without difficulty, but no information
from the interior reaches the image.  Setting $\delta$ too large
artificially shrinks the black hole shadow; setting it to zero wastes
steps near the singularity and risks floating-point overflow.  The
default $0.01\,M$ is small enough that the shadow boundary shifts by
less than a pixel at typical image resolutions, and large enough that no
ray lingers inside the horizon.
\physnote{For numerical spacetimes evolved on a finite grid, the
\texttt{OutOfGrid} constructor replaces \texttt{HitHorizon}: the ray
is terminated when it leaves the interpolatable domain rather than
when it crosses a known horizon.}

\paragraph{Escape to the celestial sphere.}

A ray that reaches large radius without being captured has escaped to
infinity.  The codebase terminates when $r > r_{\text{sky}}$, where
$r_{\text{sky}} = 500\,M$ by default.  At this radius, the
weak-field metric perturbation satisfies $2M/r_{\text{sky}} = 0.004$,
so the residual deflection beyond $r_{\text{sky}}$ is bounded by
$\order{M/r_{\text{sky}}} \approx 2 \times 10^{-3}$~radians ---
negligible compared to a pixel's angular width for any practical image
resolution.  The ray's final direction at the celestial sphere
determines which point on the sky map is sampled.
\implnote{The \texttt{TerminationConfig} record holds all six threshold
parameters: horizon radius, celestial-sphere radius, maximum affine
parameter, minimum step size, and inner and outer disc radii.  The
function \texttt{defaultTermConfig} populates these from the black hole
mass $M$.}

\paragraph{Accretion disc intersection.}

A geometrically thin accretion disc lies in the equatorial plane of the
black hole ($\theta = \pi/2$ in Boyer--Lindquist coordinates, $z = 0$
in Cartesian Kerr--Schild coordinates).  Detecting the intersection
amounts to finding where the ray crosses $z = 0$ and then checking
whether the crossing radius lies within the disc annulus
$r_{\text{in}} \leq r \leq r_{\text{out}}$.
\coderef{Geodesic.Termination}{checkDiskCrossing}

The codebase detects crossings by monitoring the sign of $z$ at each
accepted step.  If two consecutive states $(x_0^\mu, p_{\mu,0})$ and
$(x_1^\mu, p_{\mu,1})$ have $z_0$ and $z_1$ of opposite sign, the ray
crossed the equatorial plane during that step.  Since $z_0$ and $z_1$
have opposite signs, $z_0 - z_1 \neq 0$ and the fractional step
parameter
%
\begin{equation}
  s = \frac{z_0}{z_0 - z_1}
  \label{eq:disk-crossing-frac}
\end{equation}
%
is well-defined with $s \in (0, 1)$, giving $z(s) = 0$ under linear
interpolation.  This fraction measures how far through the step the
plane crossing occurred, so the crossing state is the corresponding
linear blend of the two endpoint states:
%
\begin{equation}
  x^\mu_{\text{cross}} = (1 - s)\,x^\mu_0 + s\,x^\mu_1 \,, \qquad
  p_{\mu,\text{cross}} = (1 - s)\,p_{\mu,0} + s\,p_{\mu,1} \,.
  \label{eq:disk-crossing-interp}
\end{equation}
%
The interpolated radius $r_{\text{cross}}$ is then tested against the
disc bounds.  If $r_{\text{in}} \leq r_{\text{cross}} \leq r_{\text{out}}$,
the ray terminates with \texttt{HitAccretionDisk} and the crossing state
feeds into the disc emission model; otherwise, the step is accepted
normally and integration continues.
\warnnote{The interpolated state \cref{eq:disk-crossing-interp} does not
lie exactly on the geodesic.  The linear interpolation introduces an
$\order{h^2}$ error in the crossing position, where $h$ is the local
step size --- formally below the $\order{h^5}$ local accuracy of
Dormand--Prince, but sufficient because the adaptive controller keeps
$h$ small in the strong-field regions where disc crossings occur.
The momentum interpolation is coordinate-dependent and does not respect
parallel transport; it serves only to provide approximate initial data
for the emission model.  Dense output (Hermite interpolation using the
Runge--Kutta stages) would recover $\order{h^5}$ accuracy at the
crossing point but is not implemented.}

The inner disc radius $r_{\text{in}}$ is typically set to the innermost
stable circular orbit (\cref{sec:kerr-isco}): $r_{\text{ISCO}} = 6M$ for
Schwarzschild, decreasing to $r_{\text{ISCO}} = M$ for a maximally
spinning prograde Kerr disc.  The outer radius $r_{\text{out}}$ is a
modelling choice; $r_{\text{out}} = 20M$--$50M$ captures the brightest
region of a standard thin disc.

\paragraph{Safety nets.}

Two additional criteria prevent runaway computation.  The integrator
terminates when the affine parameter exceeds a maximum
$\lambda_{\text{max}}$ (default $5000\,M$), catching rays that orbit the
photon sphere many times without escaping or plunging.  For reference,
one full orbit on the unstable null circular orbit at $r = 3M$ in
Schwarzschild corresponds to $\Delta\lambda \approx 33\,M$ (at the
$E = 1$ normalisation of \cref{sec:geo-null}), so the default allows
over a hundred near-orbits before termination --- far more than any
physically relevant ray.  The integrator also terminates when the
adaptive step size drops below $h_{\min}$ (default $10^{-12}$), which
signals dynamics the integrator cannot resolve --- typically a
near-singular region or an implementation error.  Both produce artefact
pixels that should not appear in a correctly configured render.

\paragraph{The radius function.}

All threshold comparisons use a coordinate radius $r$ computed from
the Cartesian Kerr--Schild position
$x^\mu = (t, x, y, z)$.  For the Schwarzschild metric, the radial
coordinate coincides with the Euclidean distance
$r = \sqrt{x^2 + y^2 + z^2}$
(\cref{sec:schw-kerr-schild}).  For the Kerr spacetime, the Kerr--Schild
radial coordinate $r$ --- which takes the same value as the
Boyer--Lindquist radial coordinate --- satisfies the quartic relation of
\cref{sec:kerr-bl-quartic} and must be solved at each termination check.
The codebase abstracts this through a \texttt{RadiusFn} type: a function
from position to scalar radius, defaulting to the Euclidean radius for
Schwarzschild and the quartic solver for Kerr.  This abstraction allows
the same termination logic to serve any spacetime without modification.
\xrefnote{The quartic inversion and its numerical implementation are
developed in \cref{sec:kerr-bl-quartic}.}
